\documentclass[10pt]{article}
\usepackage[margin=1in]{geometry}
\usepackage[T1]{fontenc}
\usepackage[utf8]{inputenc}
\usepackage{lmodern}
\usepackage{microtype}
\usepackage{amsmath}
\usepackage{booktabs}
\usepackage{graphicx}
\usepackage{hyperref}

\title{Optimization Follow-Up: cVAE Training Tricks Without Architecture Changes}
\author{Kirill Shumskiy, Ekaterina Petrova}
\date{\today}

\begin{document}
\maketitle

\section{Goal}
This short follow-up isolates a narrow question: can we improve the current cVAE by changing only the optimization recipe, while leaving the architecture unchanged?

The reference model is the best run from the main report, \texttt{cvae\_nb\_005}, with:
\begin{itemize}
\item latent size \(64\),
\item learning rate \(3 \cdot 10^{-4}\),
\item batch size \(64\),
\item \(\beta = 10^{-3}\),
\item the same UrbanSound8K spectrogram setup as in the main report.
\end{itemize}

\section{What Changed}
We tested a sequence of optimization-only follow-ups while keeping the architecture and data preprocessing fixed.

\paragraph{Follow-up A: warmup + free bits.}
This run added:
\begin{itemize}
\item linear KL warmup over the first 25 epochs;
\item free bits with threshold \(0.01\) per latent dimension.
\end{itemize}
Run id:
\texttt{cvae\_tricks\_20260411\_210250\_lat64\_lr0.0003\_bs64\_bkl0.001\_warm25\_fb001}.

\paragraph{Follow-up B: EMA on top of the same schedule.}
We then added exponential moving average weights with \texttt{ema\_decay=0.999},
while leaving the previous warmup/free-bits settings unchanged.
Run id:
\texttt{cvae\_ema\_20260412\_181315\_lat64\_lr0.0003\_bs64\_bkl0.001}.

\paragraph{Follow-up C: EMA + retuned schedule.}
Finally, we kept EMA but shortened and softened the KL regularization schedule:
\begin{itemize}
\item KL warmup reduced to 10 epochs;
\item free bits reduced to \(0.005\);
\item training length increased to 120 epochs;
\item patience increased to 30.
\end{itemize}
Run id:
\texttt{cvae\_ema\_tricks\_20260412\_183413\_lat64\_lr0.0003\_bs64\_bkl0.001\_warm10\_fb0005}.

\section{Results}
Table~\ref{tab:tricks_compare} compares the previous best cVAE with all three follow-up runs.

\begin{table}[h]
\centering
\caption{Best baseline cVAE vs optimization-only follow-up series.}
\label{tab:tricks_compare}
\begin{tabular}{lccccc}
\toprule
Run & Best val & Log-mel L1 & MR-STFT & FAD & Div. ratio \\
\midrule
\texttt{cvae\_nb\_005} & 0.1058 & 0.0918 & 1.8971 & \textbf{3.0198} & \textbf{0.5711} \\
\texttt{cvae\_tricks\_20260411\_210250} & 0.1106 & 0.0985 & 1.9268 & 16.0544 & 0.5136 \\
\texttt{cvae\_ema\_20260412\_181315} & 0.4291 & -- & -- & -- & -- \\
\texttt{cvae\_ema\_tricks\_20260412\_183413} & \textbf{0.1048} & \textbf{0.0909} & \textbf{1.8873} & 3.2587 & 0.5588 \\
\bottomrule
\end{tabular}
\end{table}

The results separate into three clear regimes:
\begin{itemize}
\item warmup + free bits alone hurt both reconstruction and generation badly;
\item adding EMA on top of that untuned schedule was unstable and collapsed validation quality very early;
\item EMA combined with a shorter warmup, smaller free-bits threshold, and longer training produced the best validation objective and the best reconstruction metrics of all runs.
\end{itemize}

However, the new best training objective does \emph{not} translate into the best generation score.
The strongest EMA-based run improves over \texttt{cvae\_nb\_005} on best validation loss, log-mel L1, and MR-STFT,
but its FAD is slightly worse (\(3.2587\) vs \(3.0198\)), and its diversity ratio is also slightly lower (\(0.5588\) vs \(0.5711\)).

So the updated conclusion is more nuanced than in the first follow-up:
EMA plus a better-tuned KL schedule improves reconstruction,
but the original \texttt{cvae\_nb\_005} is still slightly better as a generator under the current FAD metric.

\section{Interpretation}
This follow-up still narrows the search space in a useful way.
Two conclusions are now fairly clear.

\begin{itemize}
\item EMA is not automatically helpful in this setup; with the original warmup/free-bits schedule it was clearly harmful.
\item EMA becomes useful only after the KL schedule is retuned. In that regime it improves reconstruction, but does not yet beat the original best checkpoint on prior-matching quality.
\end{itemize}

One plausible interpretation is that the retuned EMA run learns a cleaner reconstruction manifold, while the older \texttt{cvae\_nb\_005} still matches the sampling prior slightly better under standard-normal generation.
That would explain why the new run wins on reconstruction metrics but loses on FAD.

\section{Artifacts}
The main artifacts for this follow-up series are stored under:
\begin{itemize}
\item \texttt{experiments/cvae/cvae\_tricks\_20260411\_210250\_lat64\_lr0.0003\_bs64\_bkl0.001\_warm25\_fb001/}
\item \texttt{experiments/cvae/cvae\_ema\_20260412\_181315\_lat64\_lr0.0003\_bs64\_bkl0.001/}
\item \texttt{experiments/cvae/cvae\_ema\_tricks\_20260412\_183413\_lat64\_lr0.0003\_bs64\_bkl0.001\_warm10\_fb0005/}
\item each run stores \texttt{train\_summary.json};
\item the evaluated runs also store \texttt{eval/metrics.json} and \texttt{eval/qualitative/}.
\end{itemize}

\section*{Conclusion}
The optimization-only follow-up series produced one meaningful improvement:
EMA with a retuned KL schedule gives the best validation objective and the best reconstruction metrics seen so far.
But it still does not beat the older \texttt{cvae\_nb\_005} on FAD, so the best generator in the current report remains unchanged.
The practical takeaway is that optimization tricks can help, but they are shifting the reconstruction--generation trade-off rather than solving it outright.

\end{document}
